
%%%%%%%%%%%%%%%%%%%%%%%%%%
\documentclass[12pt, a4paper]{article}
\input{systeme}
\usepackage[a4paper, top=4cm,left=2.5cm, right=2.3cm, bottom=2.7cm]{geometry}
\usepackage{amsmath, amssymb, array, multicol, enumitem}
\usepackage{pdfpages}
\usepackage{comment, graphicx}  
\usepackage{xparse}
\DeclareGraphicsExtensions{.png,.jpg}
\usepackage{pgffor}  % for loops in latex
\usepackage{ifthen}  % conditional latex compilation
\usepackage{todonotes}  % notes in bright orange
\usepackage[colorlinks=true]{hyperref}  % hyperlinks
\AtBeginDocument{% coloring of links
	\hypersetup{pdfborder={0 0 0},
		urlcolor=blue}%
}
\usepackage{fontspec}  % Hebrew fonts
\usepackage{xunicode}  % Hebrew character encoding
\usepackage{polyglossia}  % Hebrew Latex
\usepackage[absolute]{textpos}   % Proper positioning of images, e.g. BGU logo

\setdefaultlanguage{hebrew}
\setotherlanguage{english}

\newfontfamily\hebrewfont[Script=Hebrew]{Times New Roman}

\makeatletter
\def\tagform@#1{\maketag@@@{(\ignorespaces#1\unskip)}}
\makeatother


\makeatletter
\renewcommand{\SepMark}[1]{\def\@SepMark{#1}}
\makeatother
\SepMark{.}
\robustify\hebrewnumeral
\robustify\Hebrewnumeral
\robustify\Hebrewnumeralfinal

\def\urladdrname{{\itshape כתובת רשת}}
\def\emailaddrname{{\itshape כתובת דואל}}

\makeatletter
\let\orig@MakeUppercase \MakeUppercase
\makeatother
\providecommand{\MakeUppercase}[1]{#1}

\fontsize{14}{15}  

\setdefaultlanguage{hebrew}
\setotherlanguage{english}
\thispagestyle{empty}
\selectlanguage{hebrew}

\newcommand{\cc}{\mathbb{C}}
\newcommand{\CC}{\mathbb{C}}
\newcommand{\rr}{\mathbb{R}}
\newcommand{\RR}{\mathbb{R}}
\newcommand{\nn}{\mathbb{N}}
\newcommand{\zz}{\mathbb{Z}}
\newcommand{\EE}{\mathbb{E}}
\newcommand{\ZZ}{\mathbb{Z}}
\newcommand{\qq}{\mathbb{Q}}
\newcommand{\FF}{\mathbb{F}}
\newcommand{\QQ}{\mathbb{Q}}
\DeclareMathOperator{\vol}{vol}
\newcommand{\abs}[1]{\left|#1\right|}

\DeclareMathOperator{\Adj}{Adj}
\DeclareMathOperator{\deter}{det}
\DeclareMathOperator{\Img}{Im}
\DeclareMathOperator{\Ker}{Ker}
\DeclareMathOperator{\rank}{rank}
\DeclareMathOperator{\Span}{Span}
\DeclareMathOperator{\tr}{tr}


\newcommand{\calB}{\mathcal{B}}
\newcommand{\calC}{\mathcal{C}}
\newcommand{\cD}{\mathcal{D}}
\newcommand{\calE}{\mathcal{E}}


%%%% new for complex functions

\newcommand{\Var}{\mathrm{Var}}
\newcommand{\Bin}{\mathrm{Bin}}
\newcommand{\ol}{\overline}
\newcommand{\wt}{\widetilde}
\newcommand{\wh}{\widehat}
\newcommand{\re}{\mathrm{Re}}
\newcommand{\im}{\mathrm{Im}}
\newcommand{\goes}{\longrightarrow}
\newcommand{\cO}{\mathcal{O}}
\newcommand{\cR}{\mathcal{R}}
\newcommand{\Log}{\mathrm{Log}}
\newcommand{\PP}{\mathbb{P}}


%%%%%%%%%%%%%%%%%%%%%%%%%%%%%%%%%%%%


\begin{document}
\selectlanguage{hebrew}

\newcommand{\semester}{%
	סמסטר א' תשפ״ה
}
\newcommand{\course}{%
הסתברות לסטטיסטיקה}
\newcommand{\coursenumber}{9321}
\newcommand{\doctype}{%
\boxed{\text{
 דף  תרגילים 3:
 משתנים מקריים מוכרים, תוחלת ושונות}}
 \boxed{\text{
להגשה עד 28.01.25 בשעה 23:59}}
}
\newcommand{\hwknumber}{}
\newcommand{\chaptype}{%
	פרק }
\newcommand{\chapnum}{\textenglish{I}}

\newcommand{\chapname}{%
אאא } 
	
\thispagestyle{empty}
\begin{textblock*}{9cm}[0.5,0.5](90mm,3cm)
	\includegraphics[width=2in]{logo}
\end{textblock*}	

\begin{center}
	\large	
המחלקה למתמטיקה - 
	\semester  \\
	{\large \bf \course - 201.1.9921}\\
	\medskip
	{\Large  \textbf{\doctype \hwknumber}  }\\
\end{center}
\vspace{0.3cm}

%%%%%%%%%%%%%%%%%%%%%%%%%%%%%%%%%%%%

\begin{enumerate}

\item
שנה אחרי שנה פרופסור אקס מלמד את אותו הקורס ״מבוא למוטנטים״ באוניברסיטה.  
הקורס נחשב לקשה וידוע כי ההסתברות שיש סטודנט (אחד לפחות) שקיבל
$100$
בקורס היא
$0.02$,
כך זה בכל שנה באופן בלתי תלוי בשנים הקודמות בהן התקיים הקורס.
פתרו כל אחד מהסעיפים הלא קשורים הבאים.
\begin{enumerate}
\item
מה ההסתברות שרק בשנה הרביעית היה סטודנט שקיבל
$100$,
לראשונה בתולדות הקורס?
\item
מה ההסתברות שב-$6$
השנים הראשונות לא היה סטודנט שקיבל
$100$
כלל?
\item
מה ההסתברות שב-$10$ השנים הראשונות, היו בדיוק שנתיים בהן היה סטודנט שקיבל 
$100$?
\item
 בכל שנה,
 פרופ׳ מקגונגל (קולגה של פרופ׳ אקס) מלמדת במקביל אליו את הקורס ״מבוא לקסמים״.
 
ידוע כי  ההסתברות שיש לפחות סטודנט אחד שקיבל בקורס שלה
$100$
היא
$0.05$,
באופן בלתי תלוי בשנים הקודמות בהן התקיים הקורס.

עוד ידוע כי הקורסים של פרופ׳ אקס ופרופ׳ מקגונגל  בלתי תלויים אחד בשני.

שני הפרופסורים הסכימו (כל אחד/ת עבור הקורס שלו/ה) כי בשנה הראשונה  בה יש סטודנט המקבל
$100$
בקורס, הוא/היא יוצא/ת לשנת שבתון. 
\begin{enumerate}
\item[(i)]
מה ההסתברות כי שניהם יצאו לשבתון לאחר שנה בלבד?
\item[(ii)]
מה ההסתברות כי שניהם יצאו לשבתון בדיוק באותה השנה? 
\end{enumerate}
\end{enumerate}

\noindent\rule{15cm}{0.4pt}

\item
יהיו
$X\sim \Bin(4,0.3)$
ו-
$Y\sim \mathrm{Poi}(3)$
שני משתנים מקריים
ונניח כי הם
 בלתי תלויים.

חשבו את ההסתברויות הבאות:
\begin{enumerate}
\begin{multicols}{3}
\item
$\PP(X+Y=1)$.
\item
$\PP(X+Y=5)$.
\item
$\PP(X< Y)$.
\end{multicols}
\end{enumerate}

\noindent\rule{15cm}{0.4pt}

נזכיר כי
$X,Y$ מ״מ
 בלתי תלויים, אם לכל
$a,b$
 מתקיים
$.\PP(X=a, Y=b)
 =\PP(X=a)\cdot
 \PP(Y=b)$

\noindent\rule{15cm}{0.4pt}

\item
זורקים
$2$
מטבעות הוגנים ובלתי-תלויים עם 
$0$
ו-$1$
משני צידי המטבעות.
נגדיר את המשתנים המקריים הבאים:
$X_0$
מספר ה-$0$-ים שיצאו.
\quad
$X_1$
מספר ה-$1$-ים שיצאו.
\quad
$Y$
סכום המספרים שיצאו במטבעות.

$Z$
שווה
$1$
אם תוצאות המטבעות שונות,
ושווה
$0$
אם תוצאות המטבעות זהות.

בכל אחד מהסעיפים הבאים בדקו
האם 
שני המשתנים הנתונים
הם בלתי תלויים או לא, בעזרת ההגדרה.
\begin{enumerate}
\begin{multicols}{4}
\item
$X_0,X_1$.
%%%
\item
$X_0, Y$.
%%%%
\item
$X_1,Z$.
%%%%
\item
$Y,Z$.
\end{multicols}
\end{enumerate}

\noindent\rule{15cm}{0.4pt}

\newpage


\item
נדב קיבל מתנה לכבוד יום ההולדת שלו: שני כרטיסי גירוד.

בכל כרטיס גירוד ניתן לזכות ב-$100$
ש״ח בהסתברות
$0.15$,
ב-$1000$
ש״ח בהסתברות
$0.05$
או לא לזכות כלל.

נגדיר את המשתנים המקריים:

$X$
הוא מספר הכרטיסים בהם נדב זכה בסכום כספי כלשהו (חיובי).

$Y$
הוא סכום הזכייה של נדב בסך הכל משני הכרטיסים.

\begin{enumerate}
\item
מצאו את פונקציית ההסתברות של
$X$,
את התוחלת
$\EE[X]$
ואת השונות
$\Var[X]$.
%%%
\item
מצאו את פונקציית ההסתברות של
$Y$,
את התוחלת
$\EE[Y]$
ואת השונות
$\Var[Y]$.
\end{enumerate}

\noindent\rule{15cm}{0.4pt}

\item
בכל אחד מהסעיפים הלא קשורים הבאים, מצאו את פונקציית ההסתברות של
$X$,
 כאשר:
 \begin{enumerate}
%%%%%
\item
$X$
משתנה מקרי עם טווח ערכים
$\cR_X=\{0,1\}$
ונתון כי
$\EE[X]=0.25$.
%%%
\item
$X$
משתנה מקרי עם טווח ערכים
$\cR_X=\{0,1,2\}$,
ונתון כי
$\EE[X]=1.4$
וגם
$\Var[X]=0.34$.
%%%%

\end{enumerate}
\noindent\rule{15cm}{0.4pt}


\item
במחלקה לאומנות יפנית עתיקה יש
$20$
סטודנטים. כל סטודנט צריך לבחור ״קורס כללי״ אחד מתוך היצע של
$100$
קורסים כלליים שהאוניברסיטה מציעה בסמסטר.
לסטודנטים לא באמת אכפת באיזה קורס כללי יבחרו, ולכן כל סטודנט בוחר באקראי אחד מבין
$100$
הקורסים.
\begin{enumerate}
\item
מהי תוחלת מספר הקורסים הכלליים בהם לא יהיה שום סטודנט מהמחלקה לאומנות יפנית עתיקה?
\item
מהי תוחלת מספר הקורסים הכלליים בהם יהיה לפחות סטודנט אחד מהמחלקה לאומנות יפנית עתיקה?
\end{enumerate}

\noindent\rule{15cm}{0.4pt}

\item
עבור
$X\sim \Bin(n,p)$
חשבו את התוחלת
$\EE[X\cdot(X-1)]$
 בשתי דרכים שונות:
 \begin{enumerate}
 \item
 (בדרך הארוכה) בעזרת כלל לוטוס וכתיבה פורמלית של התוחלת, וזה מתחיל ככה:
 $$\EE[X(X-1)]
 =\sum_{k=0}^n
 k(k-1)
 \cdot
 \PP(X=k)$$
 ניתן להשתמש כאן בזהות בה השתמשנו בכיתה
 $p\cdot
 \binom{q}{p}
 =q\cdot \binom{q-1}{p-1}$.
 \item
 (בדרך הקצרה) בעזרת לינאריות התוחלת ושימוש בתוחלת ובשונות של
 $X$.
 \end{enumerate}

\newpage

הכוונות ורמזים:
\begin{itemize}
\item
בשאלה 1 סעיף ד׳.

אחרי כתיבה פורמלית וסימון משתנים מקריים מתאימים צריך להגיע לביטוי מהצורה
$\PP(X=Y)$,
כעת יש לחשוב על מה הם בעצם כל המקרים שמתקיים
$X=Y$,
וכך להמשיך לחשב את ההסתברות.
\item
בשאלה 2, רלוונטי לכל הסעיפים.

ראשית תארו את טווח הערכים של שני המשתנים ואז נסו לחשוב מה בעצם אומר המאורע לו רוצים לחשב הסתברות, מבחינת הערכים של כל אחד מהמשתנים, והשתמשו באי-תלות שלהם.

סעיף (ג) יותר מאתגר אך עדיין נסו לפעול לפי הכוונה זו, ואולי עדיף לעבוד עם המאורע המשלים. 
\item
בשאלה 3, רלוונטי לכל הסעיפים.

נסו לפעול ממש לפי ההגדרה של משתנים בלתי תלויים. 

זכרו כי אם החשד הוא שהמשתנים תלויים, כל שעליכם/ן להראות זה ש-
$$\PP(X=a, Y=b)\neq \PP(X=a)\cdot \PP(Y=b)$$
עבור בחירה ספציפית של
$a,b$.
לעמת זאת, אם החשד הוא שהם בלתי-תלויים עליכם/ן להראות שיש שיוויון לכל
$a,b$
שהם בטווח הערכים הרלוונטיים של המשתנים.
\item
בשאלות 4,5.

שימוש ישיר בהגדרות, לאט לאט.
\item
בשאלה 6.

בסעיף א׳, תנו סימון למשתנה המקרי הרלוונטי, חשבו איך להציג אותו בצורה אחרת כסכום של מ״מ אחרים, והשתמשו בלינאריות התוחלת.

בסעיף ב׳, תנו סימון למשתנה המקרי הרלוונטי וחשבו איך לא לעבוד קשה! הפתרון כאן צריך להיות מאוד קצר.

\item
בשאלה 7.

בסעיף א׳, יש כבר רמז והכוונה בשאלה עצמה.

בסעיף ב׳, השתמשו בלינאריות התוחלת ובהגדרת השונות. בנוסף זכרו כי תוחלת ושונת של משתנה מסוג בינומי זה משהו שאנו כבר יודעים וניתן להשתמש בזה.
\end{itemize}
\end{enumerate}


\end{document}





